\heading{数学物理方法（下）-第4次作业}


\begin{question}
已知 \((x,y,z)\) 为三维直角坐标系，判断坐标系 \((\xi,\zeta,\eta)\) 和 \((\rho',\phi',z')\) 是否为正交曲面坐标系。变换为
\[
\begin{cases}
\xi=x+A\sin kz,\\
\zeta=y-A\cos kz,\\
\eta=z,
\end{cases}
\qquad
\begin{cases}
\xi=\rho'\cos\phi',\\
\zeta=\rho'\sin\phi',\\
\eta=z'.
\end{cases}
\]
\begin{solution}
判断正交曲面坐标的基本方法，是看坐标曲线切向量
\[
\bm r_i=\frac{\partial\bm r}{\partial q_i}
\]
是否两两正交；等价地，度量张量的非对角元 \(g_{ij}=\bm r_i\cdot\bm r_j\) 是否全为零。

第一组坐标的坐标曲线切向量可由逆变换
\[
x=\xi-A\sin k\eta,\qquad y=\zeta+A\cos k\eta,\qquad z=\eta
\]
求出：
\[
\bm r_\xi=(1,0,0),\qquad
\bm r_\zeta=(0,1,0),\qquad
\bm r_\eta=(-Ak\cos k\eta,-Ak\sin k\eta,1).
\]
有
\[
\bm r_\xi\cdot\bm r_\eta=-Ak\cos k\eta,\qquad
\bm r_\zeta\cdot\bm r_\eta=-Ak\sin k\eta,
\]
一般不为零。这说明度量中存在 \(d\xi\,d\eta\) 和 \(d\zeta\,d\eta\) 的交叉项，所以 \((\xi,\zeta,\eta)\) 不是正交曲面坐标系。

第二组是柱坐标在 \((\xi,\zeta,\eta)\) 空间中的通常形式：
\[
\bm r_{\rho'}=(\cos\phi',\sin\phi',0),\quad
\bm r_{\phi'}=(-\rho'\sin\phi',\rho'\cos\phi',0),\quad
\bm r_{z'}=(0,0,1),
\]
直接计算得
\[
\bm r_{\rho'}\cdot\bm r_{\phi'}=0,\qquad
\bm r_{\rho'}\cdot\bm r_{z'}=0,\qquad
\bm r_{\phi'}\cdot\bm r_{z'}=0.
\]
其 Lamé 系数为
\[
h_{\rho'}=1,\qquad h_{\phi'}=\rho',\qquad h_{z'}=1,
\]
故 \((\rho',\phi',z')\) 是正交曲面坐标系。
\end{solution}
\end{question}

